% Part: normal-modal-logic
% Chapter: filtrations
% Section: S5-fmp

\documentclass[../../../include/open-logic-section]{subfiles}

\begin{document}

\olfileid[hi]{nml}{fil}{fmp}

\olsection{\Log{K} और \Log{S5} में परिमित प्रतिरूप गुण}

\begin{defn}
  मोडल तर्कशास्त्र के किसी तंत्र $\Sigma$ में \emph{परिमित प्रतिरूप गुण}
  है यदि जब भी कोई !!a{formula}~$!A$, $\Sigma$ के किसी प्रतिरूप के किसी
  जगत पर सत्य हो, तब $!A$, $\Sigma$ के किसी \emph{परिमित} प्रतिरूप के
  किसी जगत पर सत्य हो।
\end{defn}

\begin{prop}\ollabel{prop:K-fmp}
  \Log{K} में परिमित प्रतिरूप गुण है।
\end{prop}

\begin{proof}
  \Log{K} मान्य !!{formula}ों का समुच्चय है, अर्थात् हर प्रतिरूप \Log{K}
  का प्रतिरूप है। \olref[fil]{thm:filtrations} से, यदि
  $\mSat{M}{!A}[w]$, तो $!A$ के उप-!!{formula}ों के समुच्चय $\Gamma$ के
  माध्यम से $\mModel{M}$ के किसी भी निस्यंदन के लिए
  $\mSat{M^*}{!A}[w]$। किसी भी !!{formula} के केवल परिमिततः अनेक
  उप-!!{formula} होते हैं, अतः $\Gamma$ परिमित है।
  \olref[fin]{prop:filt-are-finite} से $\card{W^*} \le 2^n$, जहाँ $n$,
  $\Gamma$ के !!{formula}ों की संख्या है। और चूँकि \Log{K} प्रतिरूपों पर
  कोई प्रतिबंध नहीं लगाता, $\mModel{M^*}$ एक \Log{K}-प्रतिरूप है।
\end{proof}

निस्यंदनों द्वारा यह दिखाने के लिए कि किसी तर्कशास्त्र~\Log{L} में परिमित
प्रतिरूप गुण है, यह आवश्यक है कि किसी \Log{L}-प्रतिरूप का निस्यंदन स्वयं
\Log{L}-प्रतिरूप हो। प्रायः इसके लिए पर्याप्त काम करना पड़ता है और हर
निस्यंदन \Log{L}-प्रतिरूप नहीं देता। किंतु सार्विक प्रतिरूपों के लिए यह
सत्य रहता है।

\begin{prop}\ollabel{prop:univ-fin}
  $\mClass{U}$ सार्विक प्रतिरूपों का वर्ग हो
  (\olref[frd][es5]{prop:S5=univ} देखें), और
  $\mClass{U}_\mathrm{Fin}$ सभी परिमित सार्विक प्रतिरूपों का वर्ग हो। तब
  कोई भी !!{formula}~$!A$, $\mClass{U}$ में तभी मान्य है जब वह
  $\mClass{U}_\mathrm{Fin}$ में मान्य है।
\end{prop}

\begin{proof}
  परिमित सार्विक प्रतिरूप सार्विक प्रतिरूप हैं, अतः बाएँ-से-दाएँ दिशा
  तुच्छ है। दाएँ-से-बाएँ दिशा के लिए मान लें कि $!A$ किसी सार्विक
  प्रतिरूप $\mModel{M}$ के किसी जगत $w$ पर असत्य है। $\Gamma$ में $!A$
  और उसके सभी उपसूत्र रखें; स्पष्टतः $\Gamma$ परिमित है। $\mModel{M}$ का
  कोई निस्यंदन $\mModel{M^*}$ लें; तब
  \olref[fin]{prop:filt-are-finite} से $\mModel{M^*}$ परिमित है, और
  \olref[fil]{thm:filtrations} से $!A$, $\mModel{M^*}$ में $[w]$ पर
  असत्य है। केवल यह देखना शेष है कि $\mModel{M^*}$ भी सार्विक है: $u$ और
  $v$ दिए होने पर परिकल्पना से $Ruv$, और
  \olref[fil]{defn:filtration}\olref[fil]{defn:filtration-R} से
  $R^*[u][v]$ भी।
\end{proof}

\begin{cor}\ollabel{cor:S5fmp}
  \Log{S5} में परिमित प्रतिरूप गुण है।
\end{cor}

\begin{proof}
  \olref[frd][es5]{prop:S5=univ} से, यदि $!A$ किसी स्वपरावर्ती और
  यूक्लिडीय प्रतिरूप के किसी जगत पर सत्य है, तो वह किसी सार्विक प्रतिरूप
  के किसी जगत पर सत्य है। \olref{prop:univ-fin} से वह किसी परिमित सार्विक
  प्रतिरूप के किसी जगत पर सत्य है (अर्थात् $!A$ के उप-!!{formula}ों के
  समुच्चय के माध्यम से प्रतिरूप के निस्यंदन में)। प्रत्येक सार्विक प्रतिरूप
  स्वपरावर्ती और यूक्लिडीय भी है; अतः $!A$ किसी परिमित स्वपरावर्ती
  यूक्लिडीय प्रतिरूप के किसी जगत पर सत्य है।
\end{proof}

\begin{prob}
  दिखाइए कि किसी सीरियल या स्वपरावर्ती प्रतिरूप का प्रत्येक निस्यंदन भी
  क्रमशः सीरियल या स्वपरावर्ती है।
\end{prob}

\begin{prob}
  किसी सममित (संक्रमणशील, यूक्लिडीय) प्रतिरूप का असममित
  (असंक्रमणशील, गैर-यूक्लिडीय) निस्यंदन खोजिए।
\end{prob}

\end{document}
